% =============================================================================
% packages.tex — Pacotes e configurações globais para ufprdown
% Baseado no packages.tex original do template PPGInf/UFPR
% =============================================================================

% ── Codificação e fontes ─────────────────────────────────────────────────────
\usepackage[utf8]{inputenc}
\usepackage[T1]{fontenc}

\usepackage{ifthen}
\usepackage{courier}                    % Verbatim, Listings, etc

% ── Matemática ───────────────────────────────────────────────────────────────
\usepackage{amsmath}
\usepackage{amsfonts}
\usepackage{amssymb}
\usepackage{amsthm}
\usepackage{bm}                         % negrito matemático (\bm{})
\usepackage{xparse}                     % para \NewDocumentCommand (citeproc)

% Teoremas (desativado porque o bookdown usa amsthm e já define seus próprios labels
% baseados no bloco 'language:' do _bookdown.yml. Se definirmos aqui, dá erro
% de 'command \theorem already defined').

% Operadores matemáticos
\DeclareMathOperator{\E}{\mathbb{E}}
\DeclareMathOperator{\Var}{Var}
\DeclareMathOperator{\Cov}{Cov}
\DeclareMathOperator{\logit}{logit}
\DeclareMathOperator{\diag}{diag}

% ── Figuras ──────────────────────────────────────────────────────────────────
\usepackage{graphicx}
\usepackage[labelformat=simple]{subcaption}
\renewcommand\thesubfigure{(\alph{subfigure})}

% ── Tabelas ──────────────────────────────────────────────────────────────────
\usepackage{booktabs}                   % \toprule, \midrule, \bottomrule
\usepackage{longtable}                  % tabelas multi-páginas
\usepackage{multirow}                   % células mescladas
\usepackage{array}                      % extensões para tabelas

% ── Código-fonte ─────────────────────────────────────────────────────────────
\usepackage{listings}
\usepackage{xcolor}

\definecolor{codegreen}{rgb}{0,0.6,0}
\definecolor{codegray}{rgb}{0.5,0.5,0.5}
\definecolor{codepurple}{rgb}{0.58,0,0.82}
\definecolor{backcolour}{rgb}{0.97,0.97,0.95}

\lstdefinestyle{thesis-code}{
    backgroundcolor=\color{backcolour},
    commentstyle=\color{codegreen},
    keywordstyle=\color{magenta},
    numberstyle=\tiny\color{codegray},
    stringstyle=\color{codepurple},
    basicstyle=\ttfamily\footnotesize,
    breaklines=true,
    captionpos=b,
    keepspaces=true,
    numbers=left,
    numbersep=5pt,
    frame=single,
    columns=flexible,
    mathescape=true,
    inputencoding=utf8,
    extendedchars=true,
}
\lstset{style=thesis-code}

% Suporte a caracteres acentuados em listings (do template original)
\lstset{literate={á}{{\'a}}1  {ã}{{\~a}}1 {à}{{\`a}}1 {â}{{\^a}}1
                 {Á}{{\'A}}1  {Ã}{{\~A}}1 {À}{{\`A}}1 {Â}{{\^A}}1
                 {é}{{\'e}}1  {ê}{{\^e}}1 {É}{{\'E}}1  {Ê}{{\^E}}1
                 {í}{{\'\i}}1 {Í}{{\'I}}1
                 {ó}{{\'o}}1  {õ}{{\~o}}1 {ô}{{\^o}}1
                 {Ó}{{\'O}}1  {Õ}{{\~O}}1 {Ô}{{\^O}}1
                 {ú}{{\'u}}1  {Ú}{{\'U}}1
                 {ç}{{\c{c}}}1 {Ç}{{\c{C}}}1 }

% ── Algoritmos ───────────────────────────────────────────────────────────────
\usepackage{algorithm,algorithmic}
\floatname{algorithm}{Algoritmo}
\renewcommand{\algorithmiccomment}[1]{~~~// #1}

% ── Bibliografia ─────────────────────────────────────────────────────────────
% As citações são processadas pelo pandoc citeproc (pandoc >= 3.x).
% O estilo de citação é definido pelo arquivo CSL no YAML (csl: csl/apa.csl).

% ── Pacotes diversos ─────────────────────────────────────────────────────────
\usepackage{alltt}                      % verbatim melhorado
\usepackage{lipsum}                     % texto aleatório (p/ exemplos)
\usepackage[final]{pdfpages}            % inclusão de páginas em PDF
\usepackage{rotating}                   % rotação de tabelas/figuras

% ── Quadros (UFPR) ──────────────────────────────────────────────────────────
\usepackage{float}
\newfloat{quadro}{htbp}{loq}[chapter]
\floatname{quadro}{Quadro}
\newcommand{\listofquadros}{\listof{quadro}{Lista de Quadros}}

% ── Comandos personalizados ──────────────────────────────────────────────────
\newcommand{\bbeta}{\bm{\beta}}
\newcommand{\bgamma}{\bm{\gamma}}
\newcommand{\btheta}{\bm{\theta}}
\newcommand{\bSigma}{\bm{\Sigma}}
\newcommand{\TODO}[1]{\textcolor{red}{\textbf{[TODO: #1]}}}

% Atalhos úteis para centralização e fonte (notas de autor)
\newcommand{\bcenter}{\begin{center}}
\newcommand{\ecenter}{\end{center}}

% ── Anexo ───────────────────────────────────────────────% Definição de \anexo (similar ao \appendix da ppginf.cls)
\makeatletter
\newcommand{\@taganexo}{ANEXO}
\newcommand{\anexo@chapter}[1]{%
  \refstepcounter{chapter}%
  \orig@chapter*{\@taganexo\ \@Alph\c@chapter\enskip -- \enskip\MakeUppercase{#1}}%
  \addcontentsline{toc}{chapter}{\@taganexo\ \@Alph\c@chapter\enskip -- \enskip\uppercase{#1}}}%
\newcommand{\anexo}{%
  \setcounter{chapter}{0}%
  \renewcommand{\thechapter}{\@Alph\c@chapter}%
  \let\chapter\anexo@chapter%
}
\makeatother
